\section{Sharp multivariate moment asymptotics and harmonic resolution}
\label{sec:poisson-harmonic-spectrum}

The one-dimensional first-unmatched-moment theorem has the ordinary
\(r\)-th moment as its exact endpoint.  In higher dimension the quotient
topology is altered by the tail index \(d+1\) of the radial Poisson kernel.
The resulting critical moment remains explicit, and the radial structure
also permits a complete harmonic resolution of the leading tensor
coefficient.  Thus the general first-unmatched tensor is not left as an
unevaluated positive quadratic form.

We use the following entropy-transfer lemma.  It is stated separately to
make clear that the denominator need not converge uniformly to one.

\begin{lemma}[Critical two-background entropy transfer]
\label{lem:poisson-critical-two-background-transfer}
Let \(1<q\le2\), \(\kappa\ge r>0\), and \(q\kappa=2r\).  On a probability
space \((E,\Omega)\), suppose that
\[
 g_t\ge c_t>0,\qquad \int g_t\dd\Omega=1,\qquad
 c_t\longrightarrow1,\qquad \|g_t-1\|_{L^1(\Omega)}\longrightarrow0.
\]
Let \(B\in L^\infty(\Omega)\), and suppose
\[
 v_t=t^{-r}B+o(t^{-r})\quad\hbox{in }L^\infty(\Omega),
 \qquad \|e_t\|_{L^q(\Omega)}=o(t^{-\kappa}).
\]
If \(g_t+v_t+e_t\ge0\) and both \(g_t\Omega\) and
\((g_t+v_t+e_t)\Omega\) are probability measures, then
\[
 D_{\rm KL}((g_t+v_t+e_t)\Omega\|g_t\Omega)
 =\frac{t^{-2r}}2\int B^2\dd\Omega+o(t^{-2r}).
\]
\end{lemma}

\begin{proof}
Put \(\Phi(s)=(1+s)\log(1+s)-s\).  For \(1<q\le2\), convexity and the
three regions \(|z|\le1/2\), \(z<-1/2\), and \(z>1/2\) give
\begin{equation}\label{eq:poisson-critical-scalar-bregman-bound}
 0\le \Phi(u+z)-\Phi(u)-\log(1+u)z\le C_q|z|^q
\end{equation}
whenever \(|u|\le1/4\) and \(u+z\ge-1\).  Indeed, Taylor's theorem and
\(\Phi''(s)=(1+s)^{-1}\) apply in the first region, compactness applies in
the second, and \((1+z)\log(1+z)=O_q(z^q)\) applies in the third.

For \(g>0\), define \(\Psi_g(z)=g\Phi(z/g)\).  Scaling
\eqref{eq:poisson-critical-scalar-bregman-bound}, and using
\(g_t\ge1/2\) and \(\|v_t/g_t\|_\infty\le1/4\) for large \(t\), gives
\[
 0\le \Psi_{g_t}(v_t+e_t)-\Psi_{g_t}(v_t)
 -\log(1+v_t/g_t)e_t\le C_q|e_t|^q.
\]
Consequently, H\"older's inequality yields
\[
 \left|D_{\rm KL}((g_t+v_t+e_t)\Omega\|g_t\Omega)
 -\int\Psi_{g_t}(v_t)\dd\Omega\right|
 \le C\|v_t\|_{q/(q-1)}\|e_t\|_q+C\|e_t\|_q^q
 =o(t^{-2r}).
\]
Here \(r+\kappa\ge2r\) and \(q\kappa=2r\).  Uniform Taylor expansion at
zero gives
\[
 \int\Psi_{g_t}(v_t)\dd\Omega
 =\frac12\int\frac{v_t^2}{g_t}\dd\Omega+O(\|v_t\|_\infty^3).
\]
Finally,
\[
 \left|\int B^2(g_t^{-1}-1)\dd\Omega\right|
 \le2\|B\|_\infty^2\|g_t-1\|_1\longrightarrow0,
\]
which proves the assertion.
\end{proof}

For symmetric tensors, \(\odot\) denotes normalized symmetrized tensor
product, and \(I=I_d\).  A symmetric tensor is harmonic when each of its
traces vanishes.  Equivalently, the homogeneous polynomial
\(H(x)=\langle H,x^{\otimes \ell}\rangle\) is harmonic.  Every
\(T\in\operatorname{Sym}^r(\RR^d)\) has the unique Fischer decomposition
\begin{equation}\label{eq:poisson-tensor-fischer-decomposition}
 T=\sum_{j=0}^{\lfloor r/2\rfloor}I^{\odot j}\odot H_{r-2j},
 \qquad \operatorname{tr}H_{r-2j}=0.
\end{equation}
We write \((a)_k=\Gamma(a+k)/\Gamma(a)\) and use the convention
\((a)_0=1\).

\begin{theorem}[Sharp Poisson moment exponent and harmonic coefficient]
\label{thm:sharp-poisson-harmonic-moment-asymptotic}
Let \(d,r\ge1\), set
\begin{equation}\label{eq:poisson-all-order-critical-exponents}
 q_{d,r}=\min\left\{2,1+\frac{2r}{d+1}\right\},
 \qquad
 \kappa_{d,r}=\frac{2r}{q_{d,r}}
 =\max\left\{r,\frac{2r(d+1)}{d+1+2r}\right\},
\end{equation}
and let \(\nu,\eta\) be probability laws on \(\RR^d\) satisfying
\[
 \int|x|^{\kappa_{d,r}}(\nu+\eta)(\dd x)<\infty,
 \qquad
 \int x^{\otimes k}\nu(\dd x)=\int x^{\otimes k}\eta(\dd x)
 \quad(0\le k<r).
\]
Put
\[
 \Delta_r=\int x^{\otimes r}(\nu-\eta)(\dd x),
 \qquad p(y)=P_1^{(d)}(y),
\]
and define
\begin{equation}\label{eq:poisson-all-order-leading-form}
 B_{d,r,\Delta_r}(y)
 =\frac{(-1)^r}{r!p(y)}\langle\Delta_r,D^rp(y)\rangle,
 \qquad
 \mathcal Q_{d,r}(\Delta_r)
 =\frac12\int B_{d,r,\Delta_r}^2p(y)\dd y.
\end{equation}
Then
\begin{equation}\label{eq:sharp-poisson-harmonic-moment-asymptotic}
 D_{\rm KL}(P_t^{(d)}*\nu\|P_t^{(d)}*\eta)
 =\mathcal Q_{d,r}(\Delta_r)t^{-2r}+o(t^{-2r}).
\end{equation}

The quadratic form in \eqref{eq:poisson-all-order-leading-form} has the
following exact harmonic resolution.  Write
\(\Delta_r=\sum_jI^{\odot j}\odot H_{\ell_j}\) as in
\eqref{eq:poisson-tensor-fischer-decomposition}, where
\[
 \ell_j=r-2j,\qquad C_j=\frac d2+\ell_j,
\]
and put
\begin{align}
 K_{d,r,j}
 &=\frac{2^r}{r!}
 \left(\frac{d+1}{2}\right)_{\ell_j}
 \left(C_j+\frac12\right)_j(C_j)_j,
 \label{eq:poisson-harmonic-K}\\
 G_{d,r,j}(x)
 &={}_2F_1\left(-j,C_j+j+\frac12;C_j;x\right).
 \label{eq:poisson-harmonic-G}
\end{align}
Thus \(G_{d,r,j}\) is a polynomial of degree \(j\).  One has
\begin{equation}\label{eq:poisson-harmonic-quadratic-resolution}
 \mathcal Q_{d,r}(\Delta_r)
 =\sum_{j=0}^{\lfloor r/2\rfloor}
 \lambda_{d,r,j}\|H_{\ell_j}\|_{\rm HS}^2,
\end{equation}
where every weight is strictly positive and
\begin{equation}\label{eq:poisson-harmonic-weight-integral}
 \lambda_{d,r,j}
 =\frac{K_{d,r,j}^2\ell_j!}
 {2^{\ell_j+1}(d/2)_{\ell_j}B(d/2,1/2)}
 \int_0^1x^{C_j-1}(1-x)^{r-1/2}G_{d,r,j}(x)^2\dd x.
\end{equation}
In particular, the weights are finite beta sums: if
\[
 g_{j,k}=\frac{(-j)_k(C_j+j+1/2)_k}{(C_j)_k k!},
\]
then the integral in \eqref{eq:poisson-harmonic-weight-integral} equals
\begin{equation}\label{eq:poisson-harmonic-weight-beta-sum}
 \sum_{k,m=0}^jg_{j,k}g_{j,m}
 B\left(C_j+k+m,r+\frac12\right).
\end{equation}

At orders two and three, the resolution is
\begin{align}
 \lambda_{d,2,0}
 &=\frac{3(d+1)(d+3)}{4(d+5)(d+7)},
 &
 \lambda_{d,2,1}
 &=\frac{3d(d+1)(7d+9)}{4(d+3)(d+5)(d+7)},
 \label{eq:poisson-harmonic-order-two}\\
 \lambda_{d,3,0}
 &=\frac{5(d+1)(d+3)(d+5)}{4(d+7)(d+9)(d+11)},
 &
 \lambda_{d,3,1}
 &=\frac{5(d+1)(d+2)(d+3)(5d+13)}
 {4(d+5)(d+7)(d+9)(d+11)}.
 \label{eq:poisson-harmonic-order-three}
\end{align}
For example, if
\[
 \Delta_3=H_3+I\odot h,
 \qquad h=\frac3{d+2}\operatorname{tr}\Delta_3,
\]
then
\begin{equation}\label{eq:poisson-third-moment-explicit-coefficient}
 \mathcal Q_{d,3}(\Delta_3)
 =\lambda_{d,3,0}\|H_3\|_{\rm HS}^2
 +\lambda_{d,3,1}|h|^2.
\end{equation}

The exponent \(\kappa_{d,r}\) is optimal uniformly over complete moment
classes.  If \(d+1>2r\), then for every
\(r\le\rho<\kappa_{d,r}\) there are laws with finite \(\rho\)-th absolute
moments, matching tensor moments below order \(r\), and with finite nonzero
\(\Delta_r\), such that
\begin{equation}\label{eq:poisson-all-order-uniform-sharpness}
 \limsup_{t\to\infty}t^{2r}
 D_{\rm KL}(P_t^{(d)}*\nu\|P_t^{(d)}*\eta)=+\infty.
\end{equation}
For \(0<\rho<r\), the tensor \(\Delta_r\) need not be defined on the
corresponding moment class.  Hence the exponent \(r\) is also optimal when
\(d+1\le2r\).
\end{theorem}

\begin{proof}
Write \(\beta=d+1\), \(q=q_{d,r}\), \(\kappa=\kappa_{d,r}\), and
\(m=\lfloor\kappa\rfloor\).  In normalized coordinates let
\[
 F_y(z)=\frac{p(y-z)}{p(y)},\qquad
 T_m(y,z)=\sum_{k=0}^m\frac{(-1)^k}{k!p(y)}
 \langle z^{\otimes k},D^kp(y)\rangle.
\]
The two-sided Poisson tail bound gives
\begin{equation}\label{eq:poisson-critical-translate-growth}
 \|F_\cdot(z)\|_{L^q(p)}^q\asymp(1+|z|)^{\beta(q-1)}.
\end{equation}
For the upper bound, use
\(1+|y|\le(1+|z|)(1+|y-z|)\); for the lower bound, integrate over
\(|y-z|\le1\).  Repeated differentiation of
\(p(y)=c_d(1+|y|^2)^{-\beta/2}\) also shows that
\(D^kp/p\) is bounded for every fixed \(k\).  Taylor's theorem on
\(|z|\le1\), followed by \eqref{eq:poisson-critical-translate-growth} on
\(|z|>1\), therefore gives
\begin{equation}\label{eq:poisson-critical-translation-remainder}
 \|F_\cdot(z)-T_m(\cdot,z)\|_{L^q(p)}
 \le C_{d,r}|z|^\kappa.
\end{equation}
Indeed, the local power is \(m+1>\kappa\), while in the large region
\[
 \beta(1-1/q)\le\kappa,\qquad m\le\kappa.
\]

For a law \(\zeta\) satisfying the moment hypothesis, define
\[
 q_t^\zeta(y)=\int F_y(x/t)\zeta(\dd x),\qquad
 B_k^\zeta(y)=\frac{(-1)^k}{k!p(y)}
 \left\langle\int x^{\otimes k}\zeta(\dd x),D^kp(y)\right\rangle.
\]
For each fixed \(x\), the local Taylor estimate makes \(t^\kappa\) times
the remainder at \(x/t\) tend to zero, whereas
\eqref{eq:poisson-critical-translation-remainder} bounds it by
\(C|x|^\kappa\).  Minkowski's inequality and dominated convergence yield
\begin{equation}\label{eq:poisson-critical-integrated-remainder}
 q_t^\zeta=1+\sum_{k=1}^mt^{-k}B_k^\zeta+r_t^\zeta,
 \qquad \|r_t^\zeta\|_{L^q(p)}=o(t^{-\kappa}).
\end{equation}

Subtract the expansions for \(\nu\) and \(\eta\).  Moment matching gives
\[
 q_t^\nu-q_t^\eta=v_t+e_t,
 \qquad
 v_t=t^{-r}B_{d,r,\Delta_r}+o(t^{-r})
 \quad\hbox{in }L^\infty(p),
 \qquad \|e_t\|_q=o(t^{-\kappa}).
\]
The Poisson score is bounded.  Hence Jensen's inequality gives, for a
constant \(L_d\),
\[
 q_t^\eta(y)\ge \exp\{-L_d t^{-1}\textstyle\int|x|\eta(\dd x)\}
 =:c_t\longrightarrow1
\]
uniformly in \(y\).  Translation continuity of \(p\) in \(L^1\), followed
by dominated convergence in \(x\), also gives
\(\|q_t^\eta-1\|_{L^1(p)}\to0\).  Both quotients have \(p\)-integral one.
Since \(q\kappa=2r\), Lemma~\ref{lem:poisson-critical-two-background-transfer}
proves \eqref{eq:sharp-poisson-harmonic-moment-asymptotic}.

It remains to resolve the coefficient.  Let
\(a=(d+1)/2\), let \(H\) be a traceless symmetric tensor of order
\(\ell\), and write \(H(y)=\langle H,y^{\otimes\ell}\rangle\).  Contracting
the derivatives of a radial function with \(H\) eliminates every term
containing a Kronecker trace.  Thus
\begin{equation}\label{eq:harmonic-radial-derivative-identity}
 H(\partial)f(|y|^2)=2^\ell f^{(\ell)}(|y|^2)H(y).
\end{equation}
Put \(C=d/2+\ell\), \(b=a+\ell=C+1/2\),
\(s=|y|^2\), and \(x=s/(1+s)\).  Applying the radial Laplacian \(j\)
times to \eqref{eq:harmonic-radial-derivative-identity} gives
\begin{align}
 &\Delta^j\{H(y)(1+s)^{-b}\}\notag\\
 &\quad=4^j(-1)^j(b)_j(C)_jH(y)(1+s)^{-(b+j)}
 {}_2F_1(-j,C+j+1/2;C;x).
 \label{eq:iterated-harmonic-poisson-derivative}
\end{align}
For completeness, expand \((1+s)^{-b}\) at zero and use
\[
 \Delta\{H(y)s^n\}=4n(n+C-1)H(y)s^{n-1};
\]
the terminating hypergeometric expression follows after \(j\) iterations,
and the identity extends from \(0\le s<1\) to \(s\ge0\) by analyticity.

If \(T=I^{\odot j}\odot H_\ell\), where \(r=\ell+2j\), then
\(\langle T,D^rp\rangle=\Delta^jH_\ell(\partial)p\).  Equations
\eqref{eq:harmonic-radial-derivative-identity} and
\eqref{eq:iterated-harmonic-poisson-derivative} therefore give
\begin{equation}\label{eq:poisson-harmonic-mode-explicit}
 B_{d,r,T}(y)=(-1)^jK_{d,r,j}H_\ell(y)
 (1+|y|^2)^{-(\ell+j)}G_{d,r,j}
 \left(\frac{|y|^2}{1+|y|^2}\right).
\end{equation}

Different harmonic degrees are orthogonal on \(S^{d-1}\), so the mixed
terms in the Fischer decomposition vanish.  If \(\Theta\) is uniform on
\(S^{d-1}\), the traceless spherical contraction is
\[
 \EE H_\ell(\Theta)^2
 =\frac{\ell!}{2^\ell(d/2)_\ell}\|H_\ell\|_{\rm HS}^2.
\]
This follows from the spherical \(2\ell\)-th moment formula: every pairing
internal to one copy of \(H_\ell\) is a trace and vanishes, leaving exactly
\(\ell!\) cross-pairings.  By
Lemma~\ref{lem:radial-angular-contractions},
\(X=|Y|^2/(1+|Y|^2)\) has beta law with parameters \(d/2\) and \(1/2\)
when \(Y\) has density \(p\).  Substitution in
\eqref{eq:poisson-harmonic-mode-explicit} proves
\eqref{eq:poisson-harmonic-quadratic-resolution} and
\eqref{eq:poisson-harmonic-weight-integral}.  Expanding the terminating
series proves \eqref{eq:poisson-harmonic-weight-beta-sum}.  Positivity is
immediate from the integral because \(G_{d,r,j}(0)=1\).  Direct gamma
reduction gives \eqref{eq:poisson-harmonic-order-two} and
\eqref{eq:poisson-harmonic-order-three}.  The order-two formula agrees with
Theorem~\ref{thm:multidim-covariance-algebra}, since
\(H_0=(\operatorname{tr}\Delta_2)/d\).

We finish with uniform sharpness.  Suppose \(\beta>2r\) and fix
\(r\le\rho<\kappa\).  For a unit vector \(e\), put
\[
 A_R^+-A_R^-=2^{1-r}\sum_{k=0}^r(-1)^{r-k}\binom rk
 \delta_{(k+1)Re}.
\]
The positive and negative parts are probability measures, their tensor
moments agree below order \(r\), and their order-\(r\) difference is
nonzero.  Choose positive \(a_j\) with \(\sum_ja_j<\infty\), put
\(w_j=a_jR_j^{-\rho}\), and choose the radii recursively so that
\(R_{j+1}>4(r+1)R_j\).  After enlarging the first radius if necessary,
\[
 \nu=\left(1-\sum_jw_j\right)\delta_0+\sum_jw_jA_{R_j}^+,
 \qquad
 \eta=\left(1-\sum_jw_j\right)\delta_0+\sum_jw_jA_{R_j}^-
\]
are probability laws with the asserted moments.

Let \(x_j=(r+1)R_je\),
\(t_j=\delta R_jw_j^{1/\beta}<R_j/4\), and
\(E_j=B(x_j,t_j)\).  The endpoint atom of \(A_{R_j}^+\) gives
\[
 (P_{t_j}^{(d)}*\nu)(E_j)\ge c_rw_j.
\]
Every atom of \(\eta\) is at distance at least \(3R_j/4\) from \(E_j\).
The Poisson tail bound and scaling therefore give
\[
 (P_{t_j}^{(d)}*\eta)(E_j)
 \le C_{d,r}(t_j/R_j)^\beta=C_{d,r}\delta^\beta w_j.
\]
Choose \(\delta\) so that the second constant is smaller than \(c_r\).
Data processing to \(\{E_j,E_j^c\}\) and
Lemma~\ref{lem:bernoulli-kl-linear-separation} yield
\[
 D_{\rm KL}(P_{t_j}^{(d)}*\nu\|P_{t_j}^{(d)}*\eta)\ge c_{d,r}w_j.
\]
Since \(q=1+2r/\beta\) in this regime,
\[
 t_j^{2r}D_{\rm KL}(P_{t_j}^{(d)}*\nu\|P_{t_j}^{(d)}*\eta)
 \ge c_{d,r}a_j^qR_j^{2r-\rho q}.
\]
The exponent is positive precisely when \(\rho<\kappa=2r/q\).
Enlarging the radii recursively makes both \(t_j\) and the last display
tend to infinity, proving \eqref{eq:poisson-all-order-uniform-sharpness}.
For \(\rho<r\), a one-dimensional Pareto law embedded on a line has finite
\(\rho\)-th moment and infinite \(r\)-th moment.  This proves the remaining
optimality assertion.
\end{proof}
